\section{Introduction}

Antimalarial drug discovery faces a dual challenge: the Plasmodium falciparum parasite's capacity for rapid resistance evolution against single-target therapeutics, and the structural complexity of African natural products (ANPs) that have historically yielded effective but under-exploited antimalarials. Artemisinin from \textit{Artemisia annua} and quinine from \textit{Cinchona} bark exemplify this paradox—ANP-derived drugs provide clinical benefit yet their polycyclic, stereochemically rich frameworks resist translation into scalable synthetic libraries\cite{who_malaria_2024}. African ethnobotanical sources including \textit{Cryptolepis sanguinolenta}, \textit{Enantia chlorantha}, and \textit{Nauclea latifolia} contain structurally diverse alkaloids, prenylated flavonoids, and quassinoids with multitarget activity profiles\cite{betow2025}, yet computational screening pipelines default to two-dimensional molecular fingerprints that discard the three-dimensional stereochemical information critical to ANP bioactivity.

Classical cheminformatics encodes molecules as fixed-length bit vectors (e.g., ECFP4\cite{rogers2010}) or graph message-passing representations (graph neural networks, GNNs\cite{gilmer2017neural}). While effective for flat, aromatic drug-like molecules, these methods exhibit systematic blind spots for ANP structural motifs: ECFP4 radius-based hashing collapses stereoisomers into identical bit patterns, and k-hop message passing in GNNs cannot distinguish cis/trans double bonds, fused ring geometries, or axial/equatorial substituent configurations\cite{morris2019weisfeiler}. For molecules with high fractions of sp\textsuperscript{3} carbons ($F_{sp^3} > 0.5$), multiple stereocenters, and flexible polycyclic frameworks—hallmarks of ANP chemistry—this loss of three-dimensional information directly degrades predictive accuracy. Prior work in this project series demonstrated 87.5\% estimand divergence between static docking and molecular dynamics simulations for ANP-like candidates\cite{sao2026p2}, attributed to conformational flexibility and solvation effects unaccounted for by rigid 2D fingerprints. Separately, scaffold-split benchmarks showed GNN underperformance (AUC = 0.649) relative to ECFP4-Random Forest (0.689) on the same ANP-enriched dataset\cite{sao2026p5}, consistent with the 1-Weisfeiler-Lehman (1-WL) expressivity bottleneck\cite{morris2019weisfeiler} failing to generalize across complex polycyclic scaffolds.

Quantum machine learning (QML) offers a fundamentally different encoding strategy: molecules are represented as quantum states in a 2\textsuperscript{n}-dimensional Hilbert space, where entanglement can capture correlations between distant atoms—critical for long-range electrostatics, ring strain, and stereoelectronic effects in ANP structures. Quantum kernel methods\cite{havlicek2019supervised} compute similarity between quantum states via inner products, yielding kernel matrices that preserve molecular geometry without dimensionality reduction. Recent demonstrations of quantum advantage in chemical property prediction have emerged in specialized contexts\cite{boy2025,tilly2022variational}, yet systematic benchmarks on stereochemically complex natural products remain absent from the literature. Crucially, prior quantum approaches that first reduce molecules to classical fingerprints before quantum encoding\cite{sao2026p3} necessarily inherit the same stereochemical blind spots, limiting quantum enhancement to kernel structure rather than molecular representation.

Here we investigate whether quantum molecular encoding preserves African natural product structural complexity sufficiently to improve antimalarial activity prediction relative to classical baselines. We encode \num{17} candidate molecules from a curated ANP-derived panel (sourced from the upstream P1 library of \num{65856} computationally expanded structures\cite{sao2026p1}) using BondOrderMatrix representations\cite{boy2025}, where bond orders, Z/E stereochemistry, and R/S chirality are embedded directly into parameterized quantum circuits. Classification via quantum kernel support vector machines is compared against ECFP4-SVM under leave-one-out cross-validation on an intentionally small, class-imbalanced cohort designed to stress-test generalization to unseen scaffolds. This proof-of-concept Phase 1 study addresses three questions: (i) Does quantum encoding capture ANP stereochemistry that ECFP4 discards? (ii) Does this encoding translate to measurable classification advantage, or do quantum and classical kernels yield equivalent performance on real-world antimalarial screening data? (iii) For which ANP structural classes—indole alkaloids, prenylated flavonoids, polycyclic quassinoids, or simple phenolics—does quantum enhancement provide advantage versus equivalence?

Following the transparent reporting precedent established in prior quantum-inspired benchmarks\cite{sao2026p3}, we adopt an honest-negative framework: statistical equivalence or quantum underperformance are reported with identical rigor to quantum advantage, accompanied by mechanistic analysis of when and why quantum encoding succeeds or fails. This approach prioritizes falsifiable claims over promotional framing, aligning with recent calls for reproducible quantum machine learning benchmarks in drug discovery\cite{cao2022quantum}.
